\documentclass[a4paper,12pt]{article}
\usepackage{color}
\usepackage{graphicx, gensymb}
\usepackage{amsfonts, cancel, amssymb}
\usepackage{amsmath, amsthm, array, esvect, esint}
\usepackage{typearea}
\usepackage[a4paper,left=2cm,right=2cm,top=2cm,bottom=3cm,bindingoffset=5mm]{geometry}
\newcommand\numberthis{\addtocounter{equation}{1}\tag{\theequation}}
\newcommand{\bra}{}
\newcommand{\kom}{}
\newcommand{\brak}{}
\newcommand{\braket}{}
\newcommand{\intx}{}
\newcommand{\intr}{}
\newcommand{\kla}{}
\newcommand{\klam}{}
\newcommand{\klammer}{}
\newcommand{\oper}{}

\begin{document}

\title{01 Metrische und normierte Räume}
\author{Joachim Schwardt}
\date{\today}
\maketitle

\section{Vollständigkeit und Abgeschlossenheit}
Es gilt $(a)\Rightarrow(b)$:
\begin{proof}
Sei $(M, d)$ vollständig und $N$ abgeschlossen in $(M,d)$. Sei $x_n$ CF in $N$. Dann ist $x_n$ auch CF in $M$. Da $M$ vollständig ist konvergiert $x_n\rightarrow x\in M$. Da $N$ abgeschlossen ist liegt $x\in N$. Also konvergiert jede CF in N.
\end{proof}
Es gilt $(a)\Leftarrow(b)$:
\begin{proof}
Sei $(N,d)$ vollständig. Sei $x_n$ Folge in $N$, sodass $\exists x\in M:x_n\rightarrow x$. Also konvergiert $x_n$ in $M$ und ist somit insbesondere CF in $M$. Dann ist $x_n$ auch CF in $N$ und aus der Vollständigkeit folgt dann $x_n\rightarrow y\in N$. Da Grenzwerte eindeutig sind gilt $x=y$.
\end{proof}
\section{Höldersche Ungleichung}
a)
\begin{proof}
Sei $p=1$ und $q=\infty$. Dann gilt:
$$\sum_{k=1}^N|x_ky_k|\le\sum_{k=1}^N|x_k|\cdot ||y||_\infty \le ||x||_1||y||_\infty <\infty$$
Die Reihe ($N\rightarrow\infty$) konvergiert also und es gilt $xy\in l^1$.\\
Sei $||x||_p=||y||_q=1$. Dann gilt:
$$\sum_{k=1}^N|x_ky_k|\le\sum_{k=1}^N\frac{1}{p}|x_k|^p+\frac{1}{q}|y_k|^q\le\frac{1}{p}||x||_p^p+\frac{1}{q}||y||_q^q=1$$
Für allgemeine $x\in l^p,\ y\in l^q$ gilt dann:
$$\frac{\|xy\|_1}{\|x\|_p\|y\|_q}=\left\|\frac{xy}{\|x\|_p\|y\|_q}\right\|_1\le 1$$
\end{proof}
b) Minkowski Ungleichung:
\begin{proof}
Für $p=1$:
\begin{align*}
\sum_{k=1}^N|x_k+y_k|\le\sum_{k=1}^N|x_k|+|y_k|\le\|x\|_1+\|y\|_1<\infty
\end{align*}
Für $p=\infty$ (noch $\mathrm{max}_k$ auf beiden Seiten nehmen):
\begin{align*}
|x_k+y_k|\le \|x\|_\infty\|y\|_\infty
\end{align*}
Für $p\in(1,\infty)$:
\begin{align*}
\sum_{k=1}^N|x_k+y_k|^p&\le\sum_{k=1}^N\kla\left( {|x_k|+|y_k|} \right)^p\le\sum_{k=1}^N\kla\left( {2\cdot\mathrm{max}\klammer\left\{ {|x_k|,|y_k|} \right\} } \right)^p\\
&\le 2^p\sum_{k=1}^N|x_k|^p+|y_k|^p\le 2^p\kla\left( {\|x\|_p^p+\|y\|_p^p} \right)<\infty
\end{align*}
Sei $q=\frac{p}{p-1}$:
\begin{align*}
\|x+y\|_p^p&=\sum_{k=1}^\infty|x_k+y_k|^p\le\sum_{k=1}^\infty|x_k|\cdot |x_k+y_k|^{p-1}+\sum_{k=1}^\infty|y_k|\cdot |x_k+y_k|^{p-1}\\
&\kom\overset{\substack{{\text{Hölder}} \\ |}}{\le}\|x\|_p\||x+y|^{p-1}\|_q+\|y\|_p\||x+y|^{p-1}\|_q=\kla\left( {\|x\|_p+\|y\|_p} \right)\cdot\kla\left( {\sum_{k=1}^\infty|x_k+y_k|^p} \right)^\frac{1}{q}\\
&=\kla\left( {\|x\|_p+\|y\|_p} \right)\cdot \frac{\|x+y\|_p^p}{\|x+y\|_p}
\end{align*}
\end{proof}
c) 
\begin{proof}
Sei $\|x\|_p=1$. Dann gilt:
$$|x_k|=\kla\left( {|x_k|^p} \right)^{1/p}=\|x\|_p=1$$
Für $q=\infty$ gilt also schon einmal $\|x\|_q=\sup |x_k|\le 1=\|x\|_p$.\\
Sei $q<\infty$:
$$\sum |x_k|^q\le\sum|x_k|^p\le\sum_k^\infty=\|x\|_p^p=1\ \ \Rightarrow\ \ \|x\|_q\le 1=\|x\|_p$$
Sei nun $x\in l^p$ beliebig. Setze $y:=\frac{x}{\|x\|_p}$ mit $\|y\|_p=1$. Also gilt:
$$\|y\|_q\le 1\ \ \Rightarrow\ \ \|x\|_q\le\|x\|_p$$
\end{proof}
d) Banachraum
\begin{proof}
Sei $x^n$ CF in $l^p$. Für $k\in\mathbb{N}$ gilt dann:
$$|x_k^n-x_k^m|=\kla\left( {|x_k^n-x_k^m|^p} \right)^{1/p}\kom\overset{\substack{{\text{Summe über }k} \\ |}}{\le}\|x^n-x^m\|_p$$
Also ist die Folge $(x_k^n)_{n\in\mathbb{N}}$ CF in $\mathbb{K}$ und somit dort konvergent. Setze also $x_k:=\lim_{n\rightarrow\infty}x_k^n$. (Zu zeigen bleibt: $x\in l^p$ und $x^n\rightarrow x$)\\
Sei $\epsilon >0$. Dann $\exists N\in\mathbb{N}\forall n,m>N:\|x^n-x^m\|_p<\epsilon$.\\
Für $k\in\mathbb{N}$ gilt somit:
$$\kla\left( {\sum_j^k|x_j^n-x_j^m|^p} \right)^{1/p}\le\|x^n-x^m\|_p<\epsilon$$
Für $m\rightarrow\infty$ folgt dann:
$$\kla\left( {\sum_j^k|x_j^n-x_j|^p} \right)^{1/p}\le\epsilon$$
Und schließlich mit $k\rightarrow\infty$:
$$\|x^n-x\|_p\le\epsilon$$
Damit gilt $x-x^n\in l^p$ also auch $x\in l^p$ ($l^p$ ist Vektorraum), sowie $x^n\rightarrow x$.
\end{proof}

\section{Normäquivalenz}
\begin{proof}
Sei $x\in\mathbb{C}^n$. Dann $\exists j\in\klammer\left\{ {1,\dots,n} \right\}$ mit $|x_j|=\|x\|_\infty$ Damit gilt:
$$\|x\|_\infty=|x_j|\le\kla\left( {\sum_k^n|x_k|^p} \right)^{1/p}=\|x\|_p\le\kla\left( {\sum_k^n\|x\|_\infty^p} \right)^{1/p}=n^{1/p}\|x\|_\infty$$
\end{proof}
\section{Banachräume}
\begin{proof}
Die Vollständigkeit von $B(K)$ zeigt man wie die von $l^\infty$.
Wenn $C(K)$ ein abgeschlossener Teilraum von $B(K)$ ist, ist $C(K)$ auch vollständig:\\
Stetige Funktionen sind auf einer kompakten Menge beschränkt (Weierstrass), also ist $C(K)\subseteq B(K)$. \\
Sei $f_n$ eine Folge in $C(K)$ mit $f_n\rightarrow f\in B(K)$. \\
Sei $\epsilon>0$. Dann $\exists N\in\mathbb{N}:\|f_N-f\|_\infty<\epsilon$. Da $f_N$ stetig ist, $\exists\delta >0\forall s,t\in K:|t-s|<\delta\Rightarrow |f_N(t)-f_N(s)|<\epsilon$. Damit folgt:
$$|f(t)-f(s)|\le |f(t)-f_N(t)|+|f_n(t)-f_N(s)|+|f_N(s)-f(s)|\le 2\|f-f_n\|_\infty+\epsilon <3\epsilon$$ 
Also ist $f$ stetig und somit $f_n\rightarrow f\in C(K)$.
\end{proof}
\end{document}